% !TeX root = 流体力学笔记.tex
% 流体力学笔记 —— 绪论
% 流体力学的研究对象、内容与研究方法

\section{流体力学的研究对象和内容}

\subsection{流体力学是什么}

流体力学是研究流体介质的特性、状态，以及在各种力的驱动下发生的流动与质量、动量、能量输运规律的力学分支学科。

\subsection{研究对象}

流体力学的研究对象包括液态、气态与等离子态。

\section{流体力学的研究方法}

流体力学的研究方法可归纳为理论研究、实验研究、数值研究与人工智能方法四类。

\subsection{理论研究}

\begin{itemize}
	\item 建立正确可解的方程与边界条件；
	\item 方程的求解或近似解（数学上困难）。
\end{itemize}

\subsection{实验研究}

\begin{itemize}
	\item 在实验室中观察和测量流动；
	\item 认识流动、发现新现象（受环境限制，普适性差）。
\end{itemize}

\subsection{数值研究}

数值研究自 20 世纪 50 年代以后逐渐发展：
\begin{itemize}
	\item 利用数值计算模拟、预测流动，验证理论（对复杂且缺乏完善数学模型的问题不适用）；
	\item 目的：建立模型——“一切的模型都是错的”；
	\item 常用工具：
	      \begin{itemize}
		      \item 商业软件：ANSYS（Fluent、CFX）；
		      \item 开源软件：OpenFOAM（C++）；
		      \item 自编程：异构并行计算技术（CPU+GPU）。
	      \end{itemize}
\end{itemize}

\subsection{人工智能：科学研究的新范式}

人工智能将物理约束与数据学习融合，连接实验、计算与理论：
\begin{itemize}
	\item \textbf{流场重建}：融合稀疏观测与物理方程，恢复速度与压力（物理信息网络 PINN 与超分辨率重建）；
	\item \textbf{快速预测}：学习工况与流场的映射，加速重复求解（神经算子 FNO 与降阶模型）；
	\item \textbf{机理发现}：识别主导动力学，建立可解释模型；
	\item \textbf{优化控制}：根据流动反馈优化策略，服务减阻设计（强化学习与代理模型优化）。
\end{itemize}
